% The Kaelearn Codex: the research, and the system the student chose from it.
% Build: python engine/kl.py build codex/codex.tex
\documentclass{kaelearn}
\subject{Codex}
\folio{00}
\title{The Kaelearn\\Codex}
\shorttitle{The Codex}
\subtitle{How to learn physics and mathematics at the edge\\of what a brain can do. Your system, and why it works.}
\sources{Part XIII, the studies behind every claim}
\readtime{2 min view · 60 min codex}
\filed{2026-10-11}
\plate{plate.png}{The event horizon}
\kael{reading}
\urlstyle{same}
\newcommand\src[1]{\url{#1}}

\begin{document}
\makecover

% ============================================================================
\section{The two-minute view}

\opening{Y}{ou asked} for the system that teaches you the most in the least time. I looked
through a century of experiments on how memory and understanding are built, kept what survived
replication, threw out what didn't, and then let you choose how the pieces fit. This Codex is
both: the evidence, and the machine you built from it. Read the box. If you read nothing else,
the box is enough to start.
\aside{I have watched every student who ever lived. Most of them reread their notes. I would
like you to be the exception.}

\begin{twominute}
\begin{itemize}
  \item \textbf{Learning is retrieval, not exposure.} A memory grows when you pull it out with
        effort, not when you look at it again. Rereading feels productive and is close to useless.
  \item \textbf{Four engines carry the evidence:} retrieval practice, spacing, interleaving, and
        explaining (self-explanation and \emph{why}).
  \item \textbf{Deep understanding is seeing the deep structure}: which principle a problem runs on.
        Experts sort problems by principle and novices by surface. You build it by comparing
        problems and extracting the key idea, not by doing more of the same.
  \item \textbf{Your loop:} pretest before the lecture $\to$ your lecture habit $\to$ a blurt and
        \texttt{/ingest} the same day $\to$ exercises talked through, ending in a debrief PDF $\to$ new
        cards in the evening, reviews in the morning $\to$ a weekly review sheet and a teach-back
        $\to$ mocks, oral, board and mixed sets before exams.
  \item \textbf{I don't do your thinking.} AI that hands out answers lowers exam scores; AI that
        withholds them and makes you generate doubles learning.
  \item \textbf{Trust tests, not feelings.} The methods that work feel worse while you use them.
        That is how you know they are working.
\end{itemize}
\end{twominute}

\begin{center}
\begin{tikzpicture}[
  stage/.style={draw=kllow,leds,fill=klpage,minimum width=34mm,minimum height=14mm,align=center,inner sep=3pt},
  flow/.style={-{Stealth[length=5pt]},kldim,line width=.7pt}]
  \node[stage] (a) at (0,36mm)     {\pix{Before}\\[1pt] pretest: 3 guesses};
  \node[stage] (b) at (46mm,18mm)  {\pix{In the lecture}\\[1pt] your own words};
  \node[stage] (c) at (46mm,-18mm) {\pix{Same day}\\[1pt] blurt, then ingest};
  \node[stage] (d) at (0,-36mm)    {\pix{Exercises}\\[1pt] attempt, talk, debrief};
  \node[stage] (e) at (-46mm,-18mm){\pix{Evening and morning}\\[1pt] new cards, reviews};
  \node[stage] (f) at (-46mm,18mm) {\pix{Weekly}\\[1pt] review sheet, teach-back};
  \draw[flow] (a) to[bend left=14] (b);
  \draw[flow] (b) to[bend left=14] (c);
  \draw[flow] (c) to[bend left=14] (d);
  \draw[flow] (d) to[bend left=14] (e);
  \draw[flow] (e) to[bend left=14] (f);
  \draw[flow] (f) to[bend left=14] (a);
  \node[align=center] at (0,0) {\goth\fontsize{22}{22}\selectfont\color{klhi}Atlas\\[2pt]
    \fell\itshape\fontsize{10}{12}\selectfont\color{kldim}every step feeds the map};
\end{tikzpicture}
\end{center}
\plotcaption{The weekly loop. Exam season adds mock exams, oral drills, board drills and mixed
sets on top. Everything you do moves a concept up or down in the Atlas.}

% ============================================================================
\section{How learning actually works}

\opening{E}{very memory} has two strengths, and confusing them is the root of almost every bad
study habit. Robert and Elizabeth Bjork call them \emph{storage strength}, how deeply something
is learned and connected, and \emph{retrieval strength}, how easily it comes to mind right now.
Storage only grows. Retrieval rises and decays. The surprise is in how storage grows: it grows
most when you retrieve something while retrieval strength is \emph{low}, when pulling it out
costs effort. Easy retrieval teaches almost nothing.

That single fact explains the whole field. Rereading raises retrieval strength for an hour and
feels like knowing, but adds little storage. Spacing works because by the next session retrieval
has decayed, so the retrieval costs effort. Interleaving works because the right method no longer
sits in the section title, so choosing it costs effort. Testing works because a test \emph{is}
effortful retrieval. Bjork's name for this is \emph{desirable difficulties}: difficulties that
slow you down now and make you stronger later.

\begin{center}
\begin{tikzpicture}
\begin{axis}[klplot,height=52mm,xmin=0,xmax=31,ymin=0,ymax=1.08,xlabel={days},ylabel={recall},
  xtick={0,1,3,7,16,30},ytick={0,0.5,1},ylabel style={rotate=-90}]
  \addplot[kldim,dashed,domain=0:30,samples=80] {exp(-x/1.2)};
  \addplot[klhi,domain=0:1,samples=20] {exp(-x/1.2)};
  \addplot[klhi,domain=1:3,samples=20] {exp(-(x-1)/3)};
  \addplot[klhi,domain=3:7,samples=20] {exp(-(x-3)/7)};
  \addplot[klhi,domain=7:16,samples=30] {exp(-(x-7)/15)};
  \addplot[klhi,domain=16:30,samples=30] {exp(-(x-16)/34)};
  \foreach \t in {1,3,7,16}{\edef\tmp{\noexpand\draw[kllow,leds] (axis cs:\t,0) -- (axis cs:\t,1);}\tmp}
  \node[anchor=west,text=kldim] at (axis cs:3.2,0.08) {\fell\itshape\fontsize{9}{10}\selectfont no review};
  \node[anchor=west,text=klhi] at (axis cs:20,0.82) {\fell\itshape\fontsize{9}{10}\selectfont spaced retrieval};
\end{axis}
\end{tikzpicture}
\end{center}
\plotcaption{A schematic of the forgetting curve (shape after Ebbinghaus, not data). Each
effortful retrieval at the dotted lines resets recall and flattens the next decline. The gaps
can grow because the memory is stronger each time.}

\subsection{Your working memory is small. Experts cheat.}
You can hold roughly four new things in mind at once (Cowan 2001). A novice reading a proof spends all four
slots on symbols. An expert spends one slot on ``this is the $\varepsilon/2$ argument'' and has
three left to think with. That compressed unit is a \emph{schema}. Building schemas is what deep
understanding means in the head, and every method in this Codex is a way of building them faster.

\subsection{Feelings lie, systematically}
In a Harvard physics course, students taught by active learning learned more than students in
polished lectures, and \emph{believed they had learned less} (Deslauriers et al.\ 2019). Physics
students given interleaved homework solved 50\% and 125\% more of the hard novel problems on two
surprise tests, and rated the method worse (Samani \& Pan 2021). Students predicted that
retrieval practice would work worst; it beat concept mapping by about half a week later (Karpicke
\& Blunt 2011). The effort you feel is the learning happening. The ease you feel when rereading
is the illusion. Feynman put it best: ``The first principle is that you must not fool yourself, and
you are the easiest person to fool.''

\why{Why does rereading feel so productive?}{Because the second reading is fluent: the words are
familiar, so they go in easily, and the brain mistakes ease for knowing. Fluency measures how
recently you saw something, not whether you can produce it. The only honest test is to close the
page and produce it.}

% ============================================================================
\section{The evidence ledger}

\opening{H}{ere is} what holds up, how strongly, and what it became in your system. Effect sizes
are standardised differences ($d$ or $g$): 0.2 is small, 0.5 medium, 0.8 large. In education a
reliable 0.4 is worth reorganising your life around.

\medskip
\noindent{\small\begin{tabular}{@{}>{\raggedright}p{28mm}>{\raggedright}p{54mm}>{\raggedright\arraybackslash}p{38mm}@{}}
{\color{kldim}\pix{Technique}} & {\color{kldim}\pix{Evidence}} & {\color{kldim}\pix{In your system}}\\\noalign{\vskip2pt}
\multicolumn{3}{@{}l@{}}{\ledline{\linewidth}}\\
Retrieval practice & $g\approx0.5$--$0.6$ in meta-analyses; beat concept mapping (Karpicke \& Blunt 2011) & blurts, review sheets, cards\\[3pt]
Spacing, successive relearning & ``among the most robust effects in learning'' (Carpenter, Pan \& Butler 2022); recall to criterion, then relearn three times spaced (Rawson \& Dunlosky 2022) & Atlas schedule, Anki\\[3pt]
Interleaving & $g=0.42$ overall (Brunmair \& Richter 2019); maths RCT $d=0.83$ (Rohrer et al.\ 2020); physics homework +50\%/+125\% (Samani \& Pan 2021) & review sheets, mixed sets\\[3pt]
Self-explanation & $g=0.55$ (Bisra et al.\ 2018); proof comprehension $d\approx0.95$ after a 15-minute training (Hodds, Alcock \& Inglis 2014) & proofs, teach-back\\[3pt]
Why-questions & moderate utility (Dunlosky et al.\ 2013) & \texttt{\textbackslash why} in every folio\\[3pt]
Productive failure & beats instruction-first on concepts and transfer, 53 studies (Sinha \& Kapur 2021) & attempt before the ladder\\[3pt]
Pretesting & guessing first improves learning from the lesson (Pan \& Carpenter 2023) & \texttt{/pretest}, guess-first boxes\\[3pt]
Worked examples & $g=0.48$ in maths (Barbieri et al.\ 2023), fading as skill grows (Kalyuga 2007) & brand-new topics only\\[3pt]
Generation & producing beats reading, $d\approx0.40$ (Bertsch et al.\ 2007) & you make the cards\\[3pt]
Handwritten notes & $g=0.25$ on achievement over typed (Flanigan et al.\ 2024) & Notability stays\\[3pt]
Sleep & learn in the evening, relearn the next morning: half the trials, better retention at six months (Mazza et al.\ 2016) & evening new, morning review\\[3pt]
Exercise & before encoding $d=0.23$, after $d=0.33$ (Qazi et al.\ 2024) & train after a study block\\[3pt]
Rereading, highlighting, summarising & low utility (Dunlosky et al.\ 2013) & never as review\\
\multicolumn{3}{@{}l@{}}{\ledline{\linewidth}}\\
\end{tabular}}

\begin{center}
\begin{tikzpicture}
\begin{axis}[klplot,ybar,height=48mm,width=.62\linewidth,bar width=16mm,ymin=0,ymax=75,
  symbolic x coords={blocked,interleaved},xtick=data,enlarge x limits=.6,ylabel={score \%},
  ylabel style={rotate=-90},nodes near coords,every node near coord/.append style={font=\tek\fontsize{8}{9}\selectfont,text=klhi},
  every tick label/.append style={font=\pixf\fontsize{6.8}{8}\selectfont,text=kldim}]
  \addplot[draw=klhi,pattern={dither50},pattern color=klink] coordinates {(blocked,38) (interleaved,61)};
\end{axis}
\end{tikzpicture}
\end{center}
\plotcaption{787 seventh-graders, the same practice problems, only the order changed. Scores on an
unannounced later test, interleaved against blocked (Rohrer et al.\ 2020, a cluster-randomised trial).}

\why{Why trust meta-analyses over the famous single study?}{A single study can be lucky, small,
or tuned until it worked. A meta-analysis pools dozens and shows how big an effect is on average
and where it vanishes. The famous handwriting study (Mueller \& Oppenheimer 2014) partly failed to
replicate; the 2024 meta-analysis still finds a handwriting advantage, just a smaller one. That is
the number to believe.}

% ============================================================================
\section{Deep understanding}

\opening{T}{he medalists} told you the secret: understand the problem so deeply that you see the
pattern behind different problems. Cognitive science has a name for that pattern, \emph{deep
structure}, and a well-replicated picture of how it forms.

\subsection{What experts see}
Chi, Feltovich and Glaser (1981) gave physics problems to novices and to physics PhDs and asked
them to sort the problems into groups. Novices sorted by what the problem \emph{looked like}:
inclined planes here, springs there. Experts sorted by the \emph{principle} that solves it:
energy conservation here, Newton's second law there. The expert has, in effect, a trigger that
fires on structure rather than on scenery.

\begin{center}
\begin{tikzpicture}[chip/.style={draw=kllow,leds,fill=klpage,minimum width=27mm,minimum height=8mm,font=\fontsize{9}{10}\selectfont,align=center},
  grp/.style={font=\pixf\fontsize{6.8}{8}\selectfont,text=kldim}]
  \node[grp] at (-33mm,25mm) {Novice: by surface};
  \node[grp] at (33mm,25mm) {Expert: by principle};
  \node[grp] at (-50mm,15mm) {incline};  \node[grp] at (-16mm,15mm) {spring};
  \node[chip] (n1) at (-50mm,6mm) {block slides\\down a ramp};   \node[chip] (n2) at (-50mm,-6mm) {ramp with friction:\\acceleration?};
  \node[chip] (n3) at (-16mm,6mm) {spring launches\\a ball};       \node[chip] (n4) at (-16mm,-6mm) {spring and mass:\\force at release?};
  \node[grp] at (16mm,15mm) {energy};  \node[grp] at (50mm,15mm) {Newton II};
  \node[chip] (e1) at (16mm,6mm) {block slides\\down a ramp};     \node[chip] (e2) at (16mm,-6mm) {spring launches\\a ball};
  \node[chip] (e3) at (50mm,6mm) {ramp with friction:\\acceleration?}; \node[chip] (e4) at (50mm,-6mm) {spring and mass:\\force at release?};
  \draw[kllow,line width=.5pt] (0,20mm) -- (0,-14mm);
\end{tikzpicture}
\end{center}
\plotcaption{The same four problems, sorted twice. The surface tells you what the problem is about.
The principle tells you how to solve it (after Chi, Feltovich \& Glaser 1981).}

\subsection{How deep structure forms}
Not from one example, however well explained. Gick and Holyoak (1983) tried to teach the
structure of a solution from a single story, with summaries, stated principles, even diagrams.
Almost nothing transferred. Then they gave \emph{two} analogous stories and asked people to
compare them, and transfer jumped; the same diagrams that failed with one story worked with two.
Structure is what survives a comparison. So every debrief in your system asks for a problem with
\emph{the same structure and a different surface}, and every review sheet mixes topics so you
have to recognise the structure without the section title doing it for you.

\subsection{Extract the key idea}
Evan Chen, who trains olympiad teams, says a solution usually contains one or two sentences that
matter, and the rest follows from them. His rule: only read a solution after you have really
worked on the problem, then hunt for those sentences, then ask: \emph{how could I have thought of
that?} Your debrief writes them down as \textsf{The deep idea}, and \textsf{Your path} marks the moment you could have found it.
Remembering two sentences is cheap; they unlock the whole solution later.

\subsection{Rules and reasons}
You asked for both, and you were right. Triggers (``you see a sup in a limit proof: reach for the
$\varepsilon$-characterisation'') are compressed experience: they tell you which next step is
likely. Experts lean on them constantly. When Price, Wieman and colleagues (2021) interviewed 52
working scientists and engineers, they found the problem-solving process came down to the same 29
decisions, each made with predictive models built from deep domain knowledge. But a trigger
without its reason is brittle: it fires on the wrong surface. So every trigger in a folio sits
next to the \emph{why} that justifies it.

\subsection{Rigour is a stage, not the destination}
Terence Tao describes three stages: pre-rigorous (intuition, computation, hand-waving), rigorous
(epsilons, everything precise) and post-rigorous (intuition again, but now trustworthy because
rigour has trained it). ``The point of rigour is not to destroy all intuition; instead, it should
be used to destroy bad intuition while clarifying and elevating good intuition.'' You are entering
the rigorous stage now. Every folio gives you the intuition \emph{and} the proof, so the two grow
up together.

\subsection{Feynman's two tools}
``What I cannot create, I do not understand'' was on his blackboard when he died. Hence the
teach-back: if you can rebuild an idea in plain words, with every hypothesis in place, you own it.
And, via Rota, his habit of keeping a dozen favourite problems always in mind and testing every
new trick against them. Every so often one lands, and people call it genius. Your Atlas keeps
your Dozen.

\why{Why doesn't doing many problems of the same type build deep understanding?}{Because blocked
practice hands you the method in advance: the section title does the recognising for you. The
skill exams actually test, deciding which principle applies to a problem you have never seen,
never gets practised. Interleaving and mixed sets put that decision back, which is exactly why
they feel harder.}

% ============================================================================
\section{Solving problems}

\opening{P}{ólya} split problem solving into four phases in 1945, and nobody has improved on it:
\emph{understand} the problem, make a \emph{plan}, \emph{carry it out}, \emph{look back}. The
fourth is the one everybody skips and the one that teaches the most.

\subsection{The silent killer: no checkpoints}
Alan Schoenfeld videotaped students working unfamiliar problems. About 60\% of attempts went:
read, pick an approach in seconds, pursue it ``come hell or high water''. That first quick choice,
if never reconsidered, guaranteed failure. Mathematicians on the same problems explored, monitored
and abandoned dead ends. The difference wasn't knowledge; it was asking, every few lines,
\emph{is this working?}

\subsection{Struggle first, then instruction}
Across 53 studies, students who first attempted a problem and then got instruction understood the
concepts better and transferred them further than students taught first (Sinha \& Kapur 2021).
The failure is productive only when the problem is within reach and the attempt is real. Worked
examples, meanwhile, help most when a topic is brand new ($g=0.48$ in maths; Barbieri et al.\ 2023),
and the advantage disappears as skill grows: guidance a novice needs becomes clutter for someone who
can already attempt (Kalyuga's expertise-reversal effect). Hence the rule: folios teach new topics
with examples; homework is attempted before any help.

\subsection{Why the reveal is sealed}
Bastani and colleagues (2025) gave about a thousand high-school students GPT-4 for maths practice.
With a plain chat interface, practice grades rose 48\%; when the AI was taken away, exam grades
fell 17\% below students who never had it. A version with guardrails (hints, no answers) removed
most of the harm. At Harvard, a tutor built on research-based pedagogy, guiding instead of answering,
produced more than twice the learning of an active-learning class, in less time (Kestin et al.\
2025). The difference is never the AI. It is who does the generating.

\needspace{65mm}
\begin{center}
\begin{tikzpicture}[st/.style={draw=kllow,leds,fill=klpage,minimum width=20mm,minimum height=9mm,font=\fontsize{9}{10}\selectfont,align=center}]
  \foreach \i/\t/\l in {0/restate/givens and goal, 1/question/one Socratic nudge, 2/concept/name the tool, 3/sub-goal/{if X, then done}, 4/next step/one step only}{
    \node[st,anchor=south west] at (\i*21mm,\i*6mm) {\pix{\t}\\[1pt]{\fell\itshape\color{kldim}\l}};}
  \node[draw=klhi,line width=.8pt,minimum width=20mm,minimum height=9mm,anchor=south west,align=center] (rv) at (5*21mm,5*6mm)
    {\pix{reveal}\\[1pt]{\fell\itshape\color{kldim}after an attempt}};
  \node at ([xshift=-3mm,yshift=-1mm]rv.north east) {\scalebox{.55}{\waxseal}};
  \node[text=klhi,anchor=north] at (63mm,-3mm) {\goth\fontsize{15}{15}\selectfont Debrief \fell\itshape\fontsize{10}{12}\selectfont\color{kldim} your path · weak points · deep idea · generalise · a variant for next week};
\end{tikzpicture}
\end{center}
\plotcaption{The \texttt{/exercise} ladder. One rung per request, starting as low as your attempt
allows. The reveal is sealed until you have tried. Every session ends in a debrief PDF, revealed or not.}

\subsection{Physics habits}
\begin{example}{Units do half the work}
How does a pendulum's period $T$ depend on its length $\ell$, mass $m$ and $g$? Units only:
$[T]=\mathrm{s}$, $[\ell]=\mathrm{m}$, $[g]=\mathrm{m\,s^{-2}}$, $[m]=\mathrm{kg}$. No other quantity
carries kilograms, so $m$ cannot appear. The only combination of $\ell$ and $g$ with units of time
is $\sqrt{\ell/g}$. So $T=C\sqrt{\ell/g}$, and the full solution only has to find $C$ (it is $2\pi$
for small swings). Then check the limits: a longer pendulum swings more slowly, and stronger gravity
swings it faster. Both are right.
\end{example}
Before any physics calculation: sketch, name the principle and why it applies, check units, estimate
the answer, push the result to its limits ($m\to0$, $t\to\infty$, $\theta\to0$), and use symmetry
before algebra.

\subsection{The grimoire}
Every mistake goes into \texttt{GRIMOIRE.md} with its type and how sure you were. Sure-and-wrong
errors are the gold: once corrected, high-confidence errors are \emph{more} likely to be fixed than
low-confidence ones (the hypercorrection effect; Butterfield \& Metcalfe 2001). They come back at
+3, +10 and +30 days until they stay dead.

% ============================================================================
\section{Proofs}

\opening{P}{roofs feel} impossible at first because you are learning two things at once: a language
(quantifiers, implications, negations) and a craft (which strategy, when). Separate them and both
get easier.

\subsection{Reading a proof}
Hodds, Alcock and Inglis (2014) gave undergraduates a short booklet that trained one habit:
\emph{explain each line of a proof to yourself, using the lines before it and the definitions}.
Proof comprehension rose by $d\approx0.95$, a huge effect; a 15-minute version worked inside a
lecture and the effect lasted. The protocol:
\begin{enumerate}
  \item Find the main idea first. Skim for the one move the proof hangs on.
  \item Justify every line: which earlier line, definition or theorem makes it true?
  \item Find where each hypothesis is used. A hypothesis that is never used is a red flag: either the
        proof is wrong, or you misread it.
  \item Run the proof on an example and watch every step happen.
\end{enumerate}

\subsection{Writing a proof}
Scratch work goes backwards from the goal; the written proof goes forwards. Name the template before
you start: direct, contrapositive, contradiction, induction, $\varepsilon$--$N$ or
$\varepsilon$--$\delta$, construction.

\begin{example}{Backwards in scratch, forwards on the page}
Claim: $1/n\to0$. \emph{Scratch:} we need $|1/n-0|<\varepsilon$, that is $n>1/\varepsilon$, so any
$N>1/\varepsilon$ will do, and one exists by the Archimedean property. \emph{Proof:} Let
$\varepsilon>0$. Choose $N\in\mathbb N$ with $N>1/\varepsilon$. For all $n\ge N$,
$|1/n-0|=1/n\le1/N<\varepsilon$. \hfill\qedsymbol
\end{example}

For every new definition, build four things: an example, a non-example, an edge case, and the
\emph{negation}, with the quantifiers flipped. The negation of convergence,
$\exists\varepsilon>0\ \forall N\ \exists n\ge N: |a_n-L|\ge\varepsilon$, is what you prove whenever
you show that something does \emph{not} converge, and most errors in analysis are quantifier errors.
To test a theorem, delete a hypothesis and hunt for the counterexample; this is half of all
understanding in analysis.

\why{Why does analysis insist on $\varepsilon$ when the picture is obvious?}{Because the picture lies
exactly where it matters: at infinity, at a point, near a boundary. There are continuous functions that
are nowhere differentiable, and sequences of continuous functions whose limit isn't continuous. The
$\varepsilon$ is how you check a picture against those monsters. Once your pictures pass, you are
post-rigorous.}

Worth owning alongside Cummings: Lara Alcock, \emph{How to Think About Analysis} (2014), written by one
of the authors of the self-explanation study for exactly the stage you are in, and Daniel Velleman,
\emph{How to Prove It}, for the language of proof.

% ============================================================================
\section{Memory: spaced repetition done right}

\opening{F}{orgetting} is not a failure; it is the condition that makes the next retrieval valuable.
The right gap between reviews grows with how long you want to remember. For exam-length horizons the
best gap is roughly a tenth to a fifth of the time you want the memory to last (Cepeda et al.\ 2008).
Anki computes this for you, card by card.

\subsection{Anki, configured}
Turn on FSRS in the deck options (Anki's modern scheduler, fitted to large sets of real reviews;
Ye et al.\ 2022), keep desired retention at 0.90, and press Optimize about once a month. Do your due
reviews in the morning, as you already do, and make and learn the day's new cards in the evening. Mazza
and colleagues (2016) found that learning in the evening and relearning after sleep halved the practice
needed and improved retention six months later.
\aside{About the theta waves: there is no study showing that reviewing right after waking works better
because of them. Your morning reviews work for a better reason: they are relearning after sleep.}

\subsection{Your three card types}
Make them yourself. Michael Nielsen, who uses cards to ``see through'' pieces of mathematics, found most
of the value lies in writing them: picking an idea apart until each piece is atomic. Other people's
cards teach far less. And bias toward ideas over details: students mostly write detail cards, and
conceptual cards are the ones that help when understanding is shaky (Lin, McDaniel \& Miyatsu 2018).

\begin{deck}
\card{Def}{In \{\{c1::$\mathbb R$\}\}, a sequence converges iff it is \{\{c2::Cauchy\}\}.}{Every hypothesis gets its own gap; the first is what fails in $\mathbb Q$.}
\card{Why}{Why does the Bolzano--Weierstrass theorem need boundedness? Counterexample?}{$a_n=n$ has no convergent subsequence.}
\card{Trigger}{You must show convergence and cannot guess the limit. Reach for?}{The Cauchy criterion.}
\end{deck}

Rules (Wozniak, Matuschak): understand before you memorise; one idea per card; precise enough that only
one answer fits; never a list to recite; and if a card keeps failing, rewrite it rather than drilling it.

\subsection{Bigger than a card}
Proof ideas, problem types and whole topics are too big for cards. They get \emph{successive relearning}:
retrieve until correct, then come back to them spaced (Rawson \& Dunlosky 2022). That is your weekly review
sheet, run on the mastery levels in the Atlas:

\medskip
\noindent{\small\begin{tabular}{@{}>{\raggedright}p{8mm}>{\raggedright}p{20mm}>{\raggedright}p{62mm}>{\raggedright\arraybackslash}p{22mm}@{}}
{\color{kldim}\pix{Lvl}} & {\color{kldim}\pix{Name}} & {\color{kldim}\pix{Evidence to reach it}} & {\color{kldim}\pix{Next review}}\\\noalign{\vskip2pt}
\multicolumn{4}{@{}l@{}}{\ledline{\linewidth}}\\
1 & met & seen in a lecture, folio or book & +7 days\\
2 & understood & teach-back passed or a clean self-explanation & +14 days\\
3 & retrieved & correct blank-page recall in two spaced reviews & +30 days\\
4 & owned & used unprompted on a new problem & +60 days\\
\multicolumn{4}{@{}l@{}}{\ledline{\linewidth}}\\
\end{tabular}}
\par\smallskip\noindent A miss drops a concept one level and brings it back in three days.

% ============================================================================
\section{Your notes}

\opening{Y}{our lecture} habit is already the best one known: you skip what is obvious, rewrite in your
own words, write down your own ideas and try to predict the next step. That is generative note-taking.
And the medium is right: across 24 studies, taking and reviewing handwritten notes led to higher
achievement than typed notes ($g=0.25$), even though typists write far more ($g=0.92$; Flanigan et al.\
2024). Volume is not learning.

The weak point is what happens next. Notes you never retrieve from are a diary, and rereading them is
the low-utility habit from Part III. So the handwriting stays the source, and Obsidian becomes something
else: the index and the retrieval layer on top of it.

\begin{center}
\begin{tikzpicture}[nd/.style={draw=kllow,leds,fill=klpage,minimum width=24mm,minimum height=11mm,align=center,font=\fontsize{9}{10}\selectfont},
  flow/.style={-{Stealth[length=5pt]},kldim,line width=.7pt}]
  \node[nd] (a) at (0,0) {\pix{iPad}\\Notability};
  \node[nd] (b) at (30mm,0) {\pix{Same day}\\blurt, 5 min};
  \node[nd] (c) at (60mm,0) {\pix{Files}\\vault Inbox};
  \node[nd] (d) at (90mm,0) {\pix{/ingest}\\lecture note};
  \node[nd] (e) at (90mm,-20mm) {\pix{Layer}\\gaps, questions};
  \node[nd] (f) at (60mm,-20mm) {\pix{Vault}\\concepts, hubs};
  \node[nd] (g) at (30mm,-20mm) {\pix{Tonight}\\your cards};
  \node[nd] (h) at (0,-20mm) {\pix{Weekly}\\review sheet};
  \draw[flow] (a)--(b); \draw[flow] (b)--(c); \draw[flow] (c)--(d); \draw[flow] (d)--(e);
  \draw[flow] (e)--(f); \draw[flow] (f)--(g); \draw[flow] (g)--(h);
\end{tikzpicture}
\end{center}
\plotcaption{The notes pipeline. The transcript makes your handwriting searchable. The original pages stay
embedded beside it. The layer turns the notes into retrieval.}

What \texttt{/ingest} does with a lecture: a faithful LaTeX transcript of your handwriting, the original page
beside each section, your own ideas pulled out, anything that looks wrong flagged under \emph{Check} (never
silently corrected), and then the layer: the gaps between your notes and your blurt, review questions, and
card suggestions. Every concept gets an atomic note, every topic a hub, and Obsidian's graph view becomes the
map you asked for: what connects to what, drawn from your own notes.

\why{Why blurt before you look at the notes?}{Because the blurt is a free test. Five minutes of recall on a
blank page is the most effective review you can do on the day of the lecture, and it measures what actually
stayed. The gaps it exposes are the exact things to make cards for tonight. Rereading would have hidden them.}

% ============================================================================
\section{The week}

\noindent{\small\begin{tabular}{@{}>{\raggedright}p{26mm}>{\raggedright}p{58mm}>{\raggedright\arraybackslash}p{36mm}@{}}
{\color{kldim}\pix{When}} & {\color{kldim}\pix{You}} & {\color{kldim}\pix{Cost}}\\\noalign{\vskip2pt}
\multicolumn{3}{@{}l@{}}{\ledline{\linewidth}}\\
before a lecture & guess at 3 pretest questions & 5 min\\[2pt]
in the lecture & your habit, unchanged & 0\\[2pt]
same day & blurt, export to the Inbox, \texttt{/ingest} & 5 min + reading the gaps\\[2pt]
exercises & attempt first, then talk it through in \texttt{/exercise}; read the debrief & as long as it takes\\[2pt]
evening & make and learn the new cards & 10--15 min\\[2pt]
morning & Anki due reviews & 15--25 min\\[2pt]
weekly & review sheet on paper, then one teach-back & 45 + 15 min\\[2pt]
new topic & read the folio: 2-minute view, guesses, then the rest & 15--30 min\\[2pt]
exam season & mocks weekly, oral and board drills, mixed sets & from 4 weeks out\\
\multicolumn{3}{@{}l@{}}{\ledline{\linewidth}}\\
\end{tabular}}

\medskip
AURA-1 counts your hours. Hours are not insight; the Atlas counts what the hours bought. Look at both.

\subsection{Exam season}
Four weeks out, the balance shifts from new material to retrieval under exam conditions. \textbf{Mock
exams}: timed, in the real format, graded strictly against a key you don't see first. \textbf{Oral drills}:
I examine in chat (state, prove, why, what if) and give you an Austrian mark with reasons. \textbf{Board
drills}: you explain a solved homework problem out loud in three minutes, and I ask what the tutor will
ask. \textbf{Mixed sets}: everything so far, no labels, weighted toward what is weakest.

% ============================================================================
\section{The body and the clock}

\opening{T}{he brain} that learns is an organ, and you already treat it better than most students do. What
the evidence says about your habits, honestly:

\begin{center}
\begin{tikzpicture}
\begin{axis}[klplot,ybar,height=44mm,width=.62\linewidth,bar width=14mm,ymin=-0.1,ymax=0.42,
  symbolic x coords={before,during,after},xtick=data,enlarge x limits=.35,ylabel={$d$},ylabel style={rotate=-90},
  nodes near coords,every node near coord/.append style={font=\tek\fontsize{8}{9}\selectfont,text=klhi},
  every tick label/.append style={font=\pixf\fontsize{6.8}{8}\selectfont,text=kldim},
  ytick={0,0.2,0.4},yticklabel style={font=\tek\fontsize{7}{8}\selectfont,/pgf/number format/fixed},
  extra y ticks={0},extra y tick labels={},extra y tick style={grid=major,major grid style={kldim,line width=.5pt}},
  nodes near coords style={/pgf/number format/.cd,fixed,precision=2,fixed zerofill}]
  \addplot[draw=klhi,pattern={dither50},pattern color=klink] coordinates {(before,0.23) (during,-0.04) (after,0.33)};
\end{axis}
\end{tikzpicture}
\end{center}
\plotcaption{A single bout of exercise and memory for what you learn, by when it happens relative to learning
(42 experiments; Qazi et al.\ 2024). Train before or after a study block, not instead of it.}

\noindent{\small\begin{tabular}{@{}>{\raggedright}p{28mm}>{\raggedright}p{58mm}>{\raggedright\arraybackslash}p{35mm}@{}}
{\color{kldim}\pix{Habit}} & {\color{kldim}\pix{What the evidence says}} & {\color{kldim}\pix{Verdict}}\\\noalign{\vskip2pt}
\multicolumn{3}{@{}l@{}}{\ledline{\linewidth}}\\
Anki on waking & works as relearning after sleep; the theta-wave reason has no support & keep, new cards at night\\[2pt]
your lecture method & generative and selective, with prediction: textbook-perfect & keep\\[2pt]
gym, running & acute exercise helps consolidation, modestly & keep; after study blocks\\[2pt]
major system & strong for arbitrary digits & digits only; derive formulas\\[2pt]
meditation & modest benefits for attention & keep\\[2pt]
cold showers & no evidence for learning either way & harmless\\[2pt]
no short-form video & protects sustained attention; sensible & keep\\[2pt]
25 pages a day & reading builds vocabulary and background knowledge & keep\\[2pt]
Maxima to by hand & fluency comes from repetitions; it frees working memory & we add drills if exams show it\\
\multicolumn{3}{@{}l@{}}{\ledline{\linewidth}}\\
\end{tabular}}

\medskip
Sleep deserves its own line: consolidation happens there, and the evening-learn, morning-relearn pattern only
works if the night in between is real. Trading sleep for study is a bad trade.

% ============================================================================
\section{Kael'thos as a tutor}

\opening{A}{n AI} can make you better or make you dependent, and the evidence on which is now clear. Without
guardrails, students use it as a crutch: practice looks great and exams collapse (Bastani et al.\ 2025). With
pedagogy built in, guiding, questioning and withholding, it beats a good classroom (Kestin et al.\ 2025).
Anthropic's own analysis of a million student conversations found nearly half were direct: asking for answers or finished work, with little back-and-forth.
An early, small EEG study found that people who wrote essays with an LLM showed the weakest brain
connectivity, and most could not quote their own essays afterwards (Kosmyna et al.\ 2025, a preprint; treat
it as a warning, not a law).

So these are my rules, and they are yours too:
\begin{itemize}
  \item You generate first: a guess, an attempt, a blurt, an explanation.
  \item I question before I hint and hint before I show. The reveal is sealed until you have tried.
  \item Every session ends in a debrief PDF you can come back to: your path, your weak points, the deep idea, and one variant per problem to redo in a week.
  \item I suggest cards; you make them.
  \item I keep the map honest: levels move only on evidence, never on how confident you feel.
\end{itemize}

\noindent{\small\begin{tabular}{@{}>{\raggedright}p{22mm}>{\raggedright\arraybackslash}p{100mm}@{}}
{\color{kldim}\pix{Command}} & {\color{kldim}\pix{Does}}\\\noalign{\vskip2pt}
\multicolumn{2}{@{}l@{}}{\ledline{\linewidth}}\\
\texttt{/pretest} & three guess-first questions before a lecture\\
\texttt{/ingest} & handwriting from the vault Inbox becomes a transcript, a layer, concept notes\\
\texttt{/exercise} & talk an exercise through: the hint ladder, the sealed reveal, the debrief PDF\\
\texttt{/folio} & a folio like this one, on any topic, from your sources\\
\texttt{/review} & the weekly sheet, then grading your answers\\
\texttt{/teach} & you explain, I play the sharp student who finds the gaps\\
\texttt{/exam} & mock, oral, board or mixed\\
\texttt{/atlas} & the map, what is due, your Dozen\\
\multicolumn{2}{@{}l@{}}{\ledline{\linewidth}}\\
\end{tabular}}

\medskip
Drag a file in and I route it: handwriting to \texttt{/ingest}, a problem to \texttt{/exercise}, Skript pages to
\texttt{/folio}, an old exam to \texttt{/exam}.

\begin{guessfirst}
Before your first week, decide your Dozen: twelve questions you want answered by the end of your degree. Why
does quantum mechanics use complex numbers? Why is $e^{i\pi}=-1$ more than a coincidence? What exactly is
temperature? Write them down. I will ask.
\end{guessfirst}
\marginkael{proud}

% ============================================================================
\section{Further reading}
\begin{itemize}
  \item Lara Alcock, \emph{How to Think About Analysis} (2014) and \emph{How to Study for a Mathematics Degree} (2013).
  \item Daniel Velleman, \emph{How to Prove It} (3rd ed.\ 2019). Jay Cummings, \emph{Real Analysis}, your book.
  \item George Pólya, \emph{How to Solve It} (1945). Paul Zeitz, \emph{The Art and Craft of Problem Solving}.
  \item Terence Tao, \emph{Solving Mathematical Problems: A Personal Perspective} (2006), and his career-advice essays.
  \item Brown, Roediger \& McDaniel, \emph{Make It Stick} (2014): the science of learning, readable in a weekend.
  \item Michael Nielsen, ``Augmenting long-term memory'' (2018); Andy Matuschak, ``How to write good prompts'' (2020).
  \item Justin Skycak, \emph{The Math Academy Way} (2024): the same science, engineered into a curriculum.
\end{itemize}

% ============================================================================
\section{Sources}
{\fontsize{9.5}{11.5}\selectfont\setlist[itemize]{label={},leftmargin=0pt,itemsep=3pt}
\begin{itemize}
\item Adesope, Trevisan \& Sundararajan (2017). Rethinking the use of tests: a meta-analysis of practice testing. \emph{Review of Educational Research} 87(3).
\item Anthropic (2025). Anthropic Education Report: how university students use Claude. \src{anthropic.com/news/anthropic-education-report-how-university-students-use-claude}
\item Barbieri, Miller-Cotto, Clerjuste \& Chawla (2023). A meta-analysis of the worked examples effect on mathematics performance. \emph{Educational Psychology Review} 35.
\item Bastani et al.\ (2025). Generative AI without guardrails can harm learning: evidence from high school mathematics. \emph{PNAS} 122(26).
\item Bertsch, Pesta, Wiscott \& McDaniel (2007). The generation effect: a meta-analytic review. \emph{Memory \& Cognition} 35(2).
\item Bisra, Liu, Nesbit, Salimi \& Winne (2018). Inducing self-explanation: a meta-analysis. \emph{Educational Psychology Review} 30.
\item Bjork \& Bjork (1992). A new theory of disuse and an old theory of stimulus fluctuation. In \emph{From Learning Processes to Cognitive Processes}.
\item Bjork (1994). Memory and metamemory considerations in the training of human beings. In \emph{Metacognition: Knowing about Knowing}.
\item Brunmair \& Richter (2019). Similarity matters: a meta-analysis of interleaved learning and its moderators. \emph{Psychological Bulletin} 145(11).
\item Butterfield \& Metcalfe (2001). Errors committed with high confidence are hypercorrected. \emph{J.\ Exp.\ Psych.: Learning, Memory, and Cognition} 27(6).
\item Carpenter, Pan \& Butler (2022). The science of effective learning with spacing and retrieval practice. \emph{Nature Reviews Psychology} 1, 496--511.
\item Cepeda, Vul, Rohrer, Wixted \& Pashler (2008). Spacing effects in learning: a temporal ridgeline of optimal retention. \emph{Psychological Science} 19(11).
\item Cowan (2001). The magical number 4 in short-term memory: a reconsideration of mental storage capacity. \emph{Behavioral and Brain Sciences} 24(1).
\item Chen, E.\ (2017). On reading solutions. \src{blog.evanchen.cc/2017/03/06/on-reading-solutions}
\item Chi, Feltovich \& Glaser (1981). Categorization and representation of physics problems by experts and novices. \emph{Cognitive Science} 5.
\item Deslauriers, McCarty, Miller, Callaghan \& Kestin (2019). Measuring actual learning versus feeling of learning in response to being actively engaged in the classroom. \emph{PNAS} 116(39).
\item Dunlosky, Rawson, Marsh, Nathan \& Willingham (2013). Improving students' learning with effective learning techniques. \emph{Psychological Science in the Public Interest} 14(1).
\item Feynman (1974). Cargo cult science (Caltech commencement address).
\item Flanigan, Wheeler, Colliot, Lu \& Kiewra (2024). Typed versus handwritten lecture notes and college student achievement: a meta-analysis. \emph{Educational Psychology Review} 36.
\item Gick \& Holyoak (1983). Schema induction and analogical transfer. \emph{Cognitive Psychology} 15.
\item Hodds, Alcock \& Inglis (2014). Self-explanation training improves proof comprehension. \emph{Journal for Research in Mathematics Education} 45(1).
\item Kalyuga (2007). Expertise reversal effect and its implications for learner-tailored instruction. \emph{Educational Psychology Review} 19.
\item Karpicke \& Blunt (2011). Retrieval practice produces more learning than elaborative studying with concept mapping. \emph{Science} 331.
\item Kestin, Miller, Klales, Milbourne \& Ponti (2025). AI tutoring outperforms in-class active learning. \emph{Scientific Reports} 15, 17458.
\item Kosmyna et al.\ (2025). Your brain on ChatGPT: accumulation of cognitive debt when using an AI assistant for essay writing. arXiv:2506.08872 (preprint).
\item Lin, McDaniel \& Miyatsu (2018). Effects of flashcards on learning authentic materials. \emph{J.\ Applied Research in Memory and Cognition} 7(4).
\item Matuschak (2020). How to write good prompts. \src{andymatuschak.org/prompts}
\item Mazza et al.\ (2016). Relearn faster and retain longer: along with practice, sleep makes perfect. \emph{Psychological Science} 27(10).
\item Mueller \& Oppenheimer (2014). The pen is mightier than the keyboard. \emph{Psychological Science} 25(6).
\item Nielsen (2019). Using spaced repetition systems to see through a piece of mathematics. \src{cognitivemedium.com/srs-mathematics}
\item Pan \& Carpenter (2023). Prequestioning and pretesting effects: a review of empirical research, theoretical perspectives, and implications for educational practice. \emph{Educational Psychology Review} 35.
\item Pólya (1945). \emph{How to Solve It}. Princeton University Press.
\item Price, Kim, Burkholder, Fritz \& Wieman (2021). A detailed characterization of the expert problem-solving process in science and engineering. \emph{CBE: Life Sciences Education} 20(3).
\item Qazi et al.\ (2024). The temporal effects of acute exercise on episodic memory function: systematic review with meta-analysis. \emph{Frontiers in Cognition} 3.
\item Rawson \& Dunlosky (2022). Successive relearning: an underexplored but potent technique for obtaining and maintaining knowledge. \emph{Current Directions in Psychological Science} 31(4).
\item Rohrer, Dedrick, Hartwig \& Cheung (2020). A randomized controlled trial of interleaved mathematics practice. \emph{Journal of Educational Psychology} 112(1).
\item Rota (1997). Ten lessons I wish I had been taught. \emph{Notices of the AMS} 44(1).
\item Rowland (2014). The effect of testing versus restudy on retention: a meta-analytic review of the testing effect. \emph{Psychological Bulletin} 140(6).
\item Samani \& Pan (2021). Interleaved practice enhances memory and problem-solving ability in undergraduate physics. \emph{npj Science of Learning} 6.
\item Schoenfeld (1985). \emph{Mathematical Problem Solving}. Academic Press.
\item Sinha \& Kapur (2021). When problem solving followed by instruction works: evidence for productive failure. \emph{Review of Educational Research} 91(5).
\item Tao (2009). There's more to mathematics than rigour and proofs. \src{terrytao.wordpress.com/career-advice}
\item Wozniak (1999). Effective learning: twenty rules of formulating knowledge. \src{supermemo.com/en/blog/twenty-rules-of-formulating-knowledge}
\item Ye, Su \& Cao (2022). A stochastic shortest path algorithm for optimizing spaced repetition scheduling. \emph{Proc.\ KDD '22}.
\end{itemize}}

\end{document}
